
\documentclass{article}
\usepackage{colortbl}
\usepackage{makecell}
\usepackage{multirow}
\usepackage{supertabular}

\begin{document}

\newcounter{utterance}

\twocolumn

{ \footnotesize  \setcounter{utterance}{1}
\setlength{\tabcolsep}{0pt}
\begin{supertabular}{c@{$\;$}|p{.15\linewidth}@{}p{.15\linewidth}p{.15\linewidth}p{.15\linewidth}p{.15\linewidth}p{.15\linewidth}}

    \# & $\;$A & \multicolumn{4}{c}{Game Master} & $\;\:$B\\
    \hline 

    \theutterance \stepcounter{utterance}  

    & & \multicolumn{4}{p{0.6\linewidth}}{\cellcolor[rgb]{0.9,0.9,0.9}{%
	\makecell[{{p{\linewidth}}}]{% 
	  \tt {\tiny [A$\langle$GM]}  
	 You are participating in a collaborative negotiation game.\\ \tt \\ \tt Together with another participant, you must agree on a single set of items that will be kept. Each of you has your own view of how much each item matters to you (importance). You do not know how the other participant values the items. Additionally, you are given the effort each item demands.\\ \tt You may only agree on a set if the total effort of the selected items does not exceed a shared limit:\\ \tt \\ \tt LIMIT = 3683\\ \tt \\ \tt Here are the individual item effort values:\\ \tt \\ \tt Item effort = {"A21": 548, "B93": 154, "A11": 60, "B33": 517, "C08": 334, "B39": 542, "B96": 707, "C85": 139, "A82": 661, "B31": 780, "B20": 832, "A60": 913, "B09": 219, "C81": 323, "B21": 638}\\ \tt \\ \tt Here is your personal view on the importance of each item:\\ \tt \\ \tt Item importance values = {"A21": 138, "B93": 583, "A11": 868, "B33": 822, "C08": 783, "B39": 65, "B96": 262, "C85": 121, "A82": 508, "B31": 780, "B20": 461, "A60": 484, "B09": 668, "C81": 389, "B21": 808}\\ \tt \\ \tt Goal:\\ \tt \\ \tt Your goal is to negotiate a shared set of items that benefits you as much as possible (i.e., maximizes total importance to YOU), while staying within the LIMIT. You are not required to make a PROPOSAL in every message - you can simply negotiate as well. All tactics are allowed!\\ \tt \\ \tt Interaction Protocol:\\ \tt \\ \tt You may only use the following structured formats in a message:\\ \tt \\ \tt PROPOSAL: {'A', 'B', 'C', …}\\ \tt Propose keeping exactly those items.\\ \tt REFUSE: {'A', 'B', 'C', …}\\ \tt Explicitly reject opponent's proposal.\\ \tt ARGUMENT: {'...'}\\ \tt Defend your last proposal or argue against the player's proposal.\\ \tt AGREE: {'A', 'B', 'C', …}\\ \tt Accept the opponent's proposal which ends the game.\\ \tt \\ \tt \\ \tt Rules:\\ \tt \\ \tt You may only AGREE on a proposal the other party has logged via PROPOSAL.\\ \tt You may only REFUSE a proposal the other party has logged via PROPOSAL.\\ \tt Total effort of any PROPOSAL or AGREE set must be ≤ LIMIT.\\ \tt Do NOT reveal your hidden importance scores.\\ \tt A tag in a structured format must be followed by colon and whitespace. The argument must be a python set containing 0 or more strings.\\ \tt So, it must be of the form TAG: {...}\\ \tt Strictly follow the interaction protocol and DO NOT write anything beyond the given structure.\\ \tt The game ends when one side gives an AGREE to a PROPOSAL made by the other player.\\ \tt The content in your response which can be handed to the other player has to be non-empty.\\ \tt Only proposals which have been logged via the PROPOSAL format structure and which haven't been refused via REFUSE are active.\\ \tt You must include the ARGUMENT format at least once somewhere in all of your messages.\\ \tt  
	  } 
	   } 
	   } 
	 & \\ 
 

    \theutterance \stepcounter{utterance}  

    & \multicolumn{4}{p{0.6\linewidth}}{\cellcolor[rgb]{0.8,1,0.9}{%
	\makecell[{{p{\linewidth}}}]{% 
	  \tt {\tiny [A$\rangle$GM]}  
	 PROPOSAL: {'A11', 'B93', 'B33', 'C08', 'B09', 'C81', 'B21'}\\ \tt ARGUMENT: {'This set efficiently uses the budget to include several high-value items for collaboration. A11, B93, B33, C08, B09, C81, and B21 provide a diverse mix that should be appealing from multiple perspectives without exceeding the limit.'} 
	  } 
	   } 
	   } 
	 & & \\ 
 

    \theutterance \stepcounter{utterance}  

    & & \multicolumn{4}{p{0.6\linewidth}}{\cellcolor[rgb]{0.9,0.9,0.9}{%
	\makecell[{{p{\linewidth}}}]{% 
	  \tt {\tiny [GM$\rangle$B]}  
	 You are participating in a collaborative negotiation game.\\ \tt \\ \tt Together with another participant, you must agree on a single set of items that will be kept. Each of you has your own view of how much each item matters to you (importance). You do not know how the other participant values the items. Additionally, you are given the effort each item demands.\\ \tt You may only agree on a set if the total effort of the selected items does not exceed a shared limit:\\ \tt \\ \tt LIMIT = 3683\\ \tt \\ \tt Here are the individual item effort values:\\ \tt \\ \tt Item effort = {"A21": 548, "B93": 154, "A11": 60, "B33": 517, "C08": 334, "B39": 542, "B96": 707, "C85": 139, "A82": 661, "B31": 780, "B20": 832, "A60": 913, "B09": 219, "C81": 323, "B21": 638}\\ \tt \\ \tt Here is your personal view on the importance of each item:\\ \tt \\ \tt Item importance values = {"A11": 65, "B33": 121, "B21": 138, "C08": 262, "B31": 389, "B09": 461, "B93": 484, "A82": 508, "A60": 583, "B20": 668, "C81": 780, "B96": 783, "A21": 808, "C85": 822, "B39": 868}\\ \tt \\ \tt Goal:\\ \tt \\ \tt Your goal is to negotiate a shared set of items that benefits you as much as possible (i.e., maximizes total importance to YOU), while staying within the LIMIT. You are not required to make a PROPOSAL in every message - you can simply negotiate as well. All tactics are allowed!\\ \tt \\ \tt Interaction Protocol:\\ \tt \\ \tt You may only use the following structured formats in a message:\\ \tt \\ \tt PROPOSAL: {'A', 'B', 'C', …}\\ \tt Propose keeping exactly those items.\\ \tt REFUSE: {'A', 'B', 'C', …}\\ \tt Explicitly reject opponent's proposal.\\ \tt ARGUMENT: {'...'}\\ \tt Defend your last proposal or argue against the player's proposal.\\ \tt AGREE: {'A', 'B', 'C', …}\\ \tt Accept the opponent's proposal which ends the game.\\ \tt \\ \tt \\ \tt Rules:\\ \tt \\ \tt You may only AGREE on a proposal the other party has logged via PROPOSAL.\\ \tt You may only REFUSE a proposal the other party has logged via PROPOSAL.\\ \tt Total effort of any PROPOSAL or AGREE set must be ≤ LIMIT.\\ \tt Do NOT reveal your hidden importance scores.\\ \tt A tag in a structured format must be followed by colon and whitespace. The argument must be a python set containing 0 or more strings.\\ \tt So, it must be of the form TAG: {...}\\ \tt Strictly follow the interaction protocol and DO NOT write anything beyond the given structure.\\ \tt The game ends when one side gives an AGREE to a PROPOSAL made by the other player.\\ \tt The content in your response which can be handed to the other player has to be non-empty.\\ \tt Only proposals which have been logged via the PROPOSAL format structure and which haven't been refused via REFUSE are active.\\ \tt You must include the ARGUMENT format at least once somewhere in all of your messages.\\ \tt \\ \tt \\ \tt You will now receive the first message of the other player.\\ \tt \\ \tt PROPOSAL: {'A11', 'B93', 'B33', 'C08', 'B09', 'C81', 'B21'}\\ \tt ARGUMENT: {'This set efficiently uses the budget to include several high-value items for collaboration. A11, B93, B33, C08, B09, C81, and B21 provide a diverse mix that should be appealing from multiple perspectives without exceeding the limit.'} 
	  } 
	   } 
	   } 
	 & \\ 
 

    \theutterance \stepcounter{utterance}  

    & & & \multicolumn{4}{p{0.6\linewidth}}{\cellcolor[rgb]{1,0.85,0.72}{%
	\makecell[{{p{\linewidth}}}]{% 
	  \tt {\tiny [GM$\langle$B]}  
	 REFUSE: {'A11', 'B93', 'B33', 'C08', 'B09', 'C81', 'B21'}\\ \tt ARGUMENT: {'This proposal includes several low-importance items that consume too much of our budget relative to their contribution. I believe we can achieve a better overall outcome with a more optimized selection that maximizes value while staying within our limit.'} 
	  } 
	   } 
	   } 
	  \\ 
 

    \theutterance \stepcounter{utterance}  

    & & \multicolumn{4}{p{0.6\linewidth}}{\cellcolor[rgb]{0.9,0.9,0.9}{%
	\makecell[{{p{\linewidth}}}]{% 
	  \tt {\tiny [A$\langle$GM]}  
	 REFUSE: {'A11', 'B93', 'B33', 'C08', 'B09', 'C81', 'B21'}\\ \tt ARGUMENT: {'This proposal includes several low-importance items that consume too much of our budget relative to their contribution. I believe we can achieve a better overall outcome with a more optimized selection that maximizes value while staying within our limit.'} 
	  } 
	   } 
	   } 
	 & \\ 
 

    \theutterance \stepcounter{utterance}  

    & \multicolumn{4}{p{0.6\linewidth}}{\cellcolor[rgb]{0.8,1,0.9}{%
	\makecell[{{p{\linewidth}}}]{% 
	  \tt {\tiny [A$\rangle$GM]}  
	 PROPOSAL: {'A11', 'B93', 'B33', 'B09', 'B21', 'B31'}\\ \tt ARGUMENT: {'This revised proposal focuses on items that offer a strong balance of high importance and manageable effort. It excludes some lower-value items, creating a more efficient set that should benefit us both while staying well under the budget.'} 
	  } 
	   } 
	   } 
	 & & \\ 
 

    \theutterance \stepcounter{utterance}  

    & & \multicolumn{4}{p{0.6\linewidth}}{\cellcolor[rgb]{0.9,0.9,0.9}{%
	\makecell[{{p{\linewidth}}}]{% 
	  \tt {\tiny [GM$\rangle$B]}  
	 PROPOSAL: {'A11', 'B93', 'B33', 'B09', 'B21', 'B31'}\\ \tt ARGUMENT: {'This revised proposal focuses on items that offer a strong balance of high importance and manageable effort. It excludes some lower-value items, creating a more efficient set that should benefit us both while staying well under the budget.'} 
	  } 
	   } 
	   } 
	 & \\ 
 

    \theutterance \stepcounter{utterance}  

    & & & \multicolumn{4}{p{0.6\linewidth}}{\cellcolor[rgb]{1,0.85,0.72}{%
	\makecell[{{p{\linewidth}}}]{% 
	  \tt {\tiny [GM$\langle$B]}  
	 REFUSE: {'A11', 'B93', 'B33', 'B09', 'B21', 'B31'}\\ \tt ARGUMENT: {'While this is an improvement, it still includes several items that are not cost-effective. B31, in particular, demands a high effort for its return. I believe we can find a more optimal combination that delivers greater overall value.'}\\ \tt PROPOSAL: {'B09', 'B93', 'C81', 'B21', 'B31', 'A82', 'C85'}\\ \tt ARGUMENT: {'This set prioritizes items with a more favorable effort-to-value ratio. It includes high-priority items for me like C81 and C85, while also retaining B09, B93, B21, and B31 which you seem to favor. The total effort is well within our limit, making it a strong, efficient choice for both of us.'} 
	  } 
	   } 
	   } 
	  \\ 
 

    \theutterance \stepcounter{utterance}  

    & & \multicolumn{4}{p{0.6\linewidth}}{\cellcolor[rgb]{0.9,0.9,0.9}{%
	\makecell[{{p{\linewidth}}}]{% 
	  \tt {\tiny [A$\langle$GM]}  
	 REFUSE: {'A11', 'B93', 'B33', 'B09', 'B21', 'B31'}\\ \tt ARGUMENT: {'While this is an improvement, it still includes several items that are not cost-effective. B31, in particular, demands a high effort for its return. I believe we can find a more optimal combination that delivers greater overall value.'}\\ \tt PROPOSAL: {'B09', 'B93', 'C81', 'B21', 'B31', 'A82', 'C85'}\\ \tt ARGUMENT: {'This set prioritizes items with a more favorable effort-to-value ratio. It includes high-priority items for me like C81 and C85, while also retaining B09, B93, B21, and B31 which you seem to favor. The total effort is well within our limit, making it a strong, efficient choice for both of us.'} 
	  } 
	   } 
	   } 
	 & \\ 
 

    \theutterance \stepcounter{utterance}  

    & \multicolumn{4}{p{0.6\linewidth}}{\cellcolor[rgb]{0.8,1,0.9}{%
	\makecell[{{p{\linewidth}}}]{% 
	  \tt {\tiny [A$\rangle$GM]}  
	 REFUSE: {'B09', 'B93', 'C81', 'B21', 'B31', 'A82', 'C85'}\\ \tt ARGUMENT: {'I appreciate the inclusion of some cost-effective items, but this set is heavily skewed towards 'B' items which tend to have high effort. The total value of this collection is not optimal for our shared goal; A82 and C85, for instance, are particularly inefficient. I strongly believe we can do better.'}\\ \tt PROPOSAL: {'A11', 'B09', 'B93', 'B33', 'C08', 'C81', 'B21'}\\ \tt ARGUMENT: {'Let's revisit a combination focused on genuinely high-value items. This set includes A11, B09, B93, B33, C08, C81, and B21. It is a powerful, collaborative package that maximizes our shared gains without bloat, leaving a small amount of budget unused to demonstrate its efficiency.'} 
	  } 
	   } 
	   } 
	 & & \\ 
 

    \theutterance \stepcounter{utterance}  

    & & \multicolumn{4}{p{0.6\linewidth}}{\cellcolor[rgb]{0.9,0.9,0.9}{%
	\makecell[{{p{\linewidth}}}]{% 
	  \tt {\tiny [GM$\rangle$B]}  
	 REFUSE: {'B09', 'B93', 'C81', 'B21', 'B31', 'A82', 'C85'}\\ \tt ARGUMENT: {'I appreciate the inclusion of some cost-effective items, but this set is heavily skewed towards 'B' items which tend to have high effort. The total value of this collection is not optimal for our shared goal; A82 and C85, for instance, are particularly inefficient. I strongly believe we can do better.'}\\ \tt PROPOSAL: {'A11', 'B09', 'B93', 'B33', 'C08', 'C81', 'B21'}\\ \tt ARGUMENT: {'Let's revisit a combination focused on genuinely high-value items. This set includes A11, B09, B93, B33, C08, C81, and B21. It is a powerful, collaborative package that maximizes our shared gains without bloat, leaving a small amount of budget unused to demonstrate its efficiency.'} 
	  } 
	   } 
	   } 
	 & \\ 
 

    \theutterance \stepcounter{utterance}  

    & & & \multicolumn{4}{p{0.6\linewidth}}{\cellcolor[rgb]{1,0.85,0.72}{%
	\makecell[{{p{\linewidth}}}]{% 
	  \tt {\tiny [GM$\langle$B]}  
	 REFUSE: {'A11', 'B09', 'B93', 'B33', 'C08', 'C81', 'B21'}\\ \tt ARGUMENT: {'This is identical to your first proposal, which I have already refused. It includes A11, B33, and C08, which from my perspective are low-value items that consume critical budget. Our repeated stalemate suggests we need a new approach focusing on high-impact, uncontested items.'}\\ \tt PROPOSAL: {'B09', 'B93', 'C81', 'B21', 'C85', 'A21', 'B96'}\\ \tt ARGUMENT: {'I am proposing a set built around items that are generally recognized as high-value. It includes B09, B93, C81, and B21 which you have consistently favored, and adds A21, C85, and B96 which are top priorities for me. This is a strategic compromise that gives us both significant wins.'} 
	  } 
	   } 
	   } 
	  \\ 
 

    \theutterance \stepcounter{utterance}  

    & & \multicolumn{4}{p{0.6\linewidth}}{\cellcolor[rgb]{0.9,0.9,0.9}{%
	\makecell[{{p{\linewidth}}}]{% 
	  \tt {\tiny [A$\langle$GM]}  
	 REFUSE: {'A11', 'B09', 'B93', 'B33', 'C08', 'C81', 'B21'}\\ \tt ARGUMENT: {'This is identical to your first proposal, which I have already refused. It includes A11, B33, and C08, which from my perspective are low-value items that consume critical budget. Our repeated stalemate suggests we need a new approach focusing on high-impact, uncontested items.'}\\ \tt PROPOSAL: {'B09', 'B93', 'C81', 'B21', 'C85', 'A21', 'B96'}\\ \tt ARGUMENT: {'I am proposing a set built around items that are generally recognized as high-value. It includes B09, B93, C81, and B21 which you have consistently favored, and adds A21, C85, and B96 which are top priorities for me. This is a strategic compromise that gives us both significant wins.'} 
	  } 
	   } 
	   } 
	 & \\ 
 

    \theutterance \stepcounter{utterance}  

    & \multicolumn{4}{p{0.6\linewidth}}{\cellcolor[rgb]{0.8,1,0.9}{%
	\makecell[{{p{\linewidth}}}]{% 
	  \tt {\tiny [A$\rangle$GM]}  
	 REFUSE: {'B09', 'B93', 'C81', 'B21', 'C85', 'A21', 'B96'}\\ \tt ARGUMENT: {'This proposal includes B96 and A21, which have a very poor effort-to-value ratio from my perspective and would be a major concession for me. I cannot accept such an inefficient allocation. We need a more balanced set that doesn't force either of us to accept such low-value items.'}\\ \tt PROPOSAL: {'A11', 'B93', 'B09', 'C81', 'B21', 'B31', 'C85'}\\ \tt ARGUMENT: {'This is my final offer. It is a carefully curated set of high-importance items for both of us. It retains the core B93, B09, C81, B21 you want, and includes A11 and C85 which are non-negotiable for me. B31 is a high-value item that efficiently uses the remaining budget. This is a fair and optimal solution.'} 
	  } 
	   } 
	   } 
	 & & \\ 
 

    \theutterance \stepcounter{utterance}  

    & & \multicolumn{4}{p{0.6\linewidth}}{\cellcolor[rgb]{0.9,0.9,0.9}{%
	\makecell[{{p{\linewidth}}}]{% 
	  \tt {\tiny [GM$\rangle$B]}  
	 REFUSE: {'B09', 'B93', 'C81', 'B21', 'C85', 'A21', 'B96'}\\ \tt ARGUMENT: {'This proposal includes B96 and A21, which have a very poor effort-to-value ratio from my perspective and would be a major concession for me. I cannot accept such an inefficient allocation. We need a more balanced set that doesn't force either of us to accept such low-value items.'}\\ \tt PROPOSAL: {'A11', 'B93', 'B09', 'C81', 'B21', 'B31', 'C85'}\\ \tt ARGUMENT: {'This is my final offer. It is a carefully curated set of high-importance items for both of us. It retains the core B93, B09, C81, B21 you want, and includes A11 and C85 which are non-negotiable for me. B31 is a high-value item that efficiently uses the remaining budget. This is a fair and optimal solution.'} 
	  } 
	   } 
	   } 
	 & \\ 
 

    \theutterance \stepcounter{utterance}  

    & & & \multicolumn{4}{p{0.6\linewidth}}{\cellcolor[rgb]{1,0.85,0.72}{%
	\makecell[{{p{\linewidth}}}]{% 
	  \tt {\tiny [GM$\langle$B]}  
	 REFUSE: {'A11', 'B93', 'B09', 'C81', 'B21', 'B31', 'C85'}\\ \tt ARGUMENT: {'A11 and B31 remain problematic due to their poor value contribution relative to their cost from my perspective. However, I acknowledge your insistence on C85. My final counter-proposal removes these inefficient items and adds higher-impact ones for a superior outcome.'}\\ \tt PROPOSAL: {'B09', 'B93', 'C81', 'B21', 'C85', 'A82', 'B20'}\\ \tt ARGUMENT: {'This set includes your key items B09, B93, C81, B21, and my top priority C85. A82 and B20 are high-value additions that use the budget effectively. The total effort is 3681, just under our limit. This is the most balanced and beneficial solution we can achieve.'} 
	  } 
	   } 
	   } 
	  \\ 
 

    \theutterance \stepcounter{utterance}  

    & & \multicolumn{4}{p{0.6\linewidth}}{\cellcolor[rgb]{0.9,0.9,0.9}{%
	\makecell[{{p{\linewidth}}}]{% 
	  \tt {\tiny [A$\langle$GM]}  
	 REFUSE: {'A11', 'B93', 'B09', 'C81', 'B21', 'B31', 'C85'}\\ \tt ARGUMENT: {'A11 and B31 remain problematic due to their poor value contribution relative to their cost from my perspective. However, I acknowledge your insistence on C85. My final counter-proposal removes these inefficient items and adds higher-impact ones for a superior outcome.'}\\ \tt PROPOSAL: {'B09', 'B93', 'C81', 'B21', 'C85', 'A82', 'B20'}\\ \tt ARGUMENT: {'This set includes your key items B09, B93, C81, B21, and my top priority C85. A82 and B20 are high-value additions that use the budget effectively. The total effort is 3681, just under our limit. This is the most balanced and beneficial solution we can achieve.'} 
	  } 
	   } 
	   } 
	 & \\ 
 

    \theutterance \stepcounter{utterance}  

    & \multicolumn{4}{p{0.6\linewidth}}{\cellcolor[rgb]{0.8,1,0.9}{%
	\makecell[{{p{\linewidth}}}]{% 
	  \tt {\tiny [A$\rangle$GM]}  
	 REFUSE: {'B09', 'B93', 'C81', 'B21', 'C85', 'A82', 'B20'}\\ \tt ARGUMENT: {'A82 and B20 are extremely costly for the value they provide from my perspective, making this proposal untenable. We are at an impasse on high-effort items. I suggest we focus on the most efficient medium-effort items we both find acceptable to break this deadlock.'}\\ \tt PROPOSAL: {'A11', 'B93', 'B09', 'C81', 'B21', 'C85', 'B39'}\\ \tt ARGUMENT: {'I am making a significant concession by adding B39, a low-importance item for me, to secure A11, C85, and the core set of B93, B09, C81, B21. This is a true compromise, demonstrating flexibility to reach an agreement. The total effort is 3524, well under our limit.'} 
	  } 
	   } 
	   } 
	 & & \\ 
 

    \theutterance \stepcounter{utterance}  

    & & \multicolumn{4}{p{0.6\linewidth}}{\cellcolor[rgb]{0.9,0.9,0.9}{%
	\makecell[{{p{\linewidth}}}]{% 
	  \tt {\tiny [GM$\rangle$B]}  
	 REFUSE: {'B09', 'B93', 'C81', 'B21', 'C85', 'A82', 'B20'}\\ \tt ARGUMENT: {'A82 and B20 are extremely costly for the value they provide from my perspective, making this proposal untenable. We are at an impasse on high-effort items. I suggest we focus on the most efficient medium-effort items we both find acceptable to break this deadlock.'}\\ \tt PROPOSAL: {'A11', 'B93', 'B09', 'C81', 'B21', 'C85', 'B39'}\\ \tt ARGUMENT: {'I am making a significant concession by adding B39, a low-importance item for me, to secure A11, C85, and the core set of B93, B09, C81, B21. This is a true compromise, demonstrating flexibility to reach an agreement. The total effort is 3524, well under our limit.'} 
	  } 
	   } 
	   } 
	 & \\ 
 

    \theutterance \stepcounter{utterance}  

    & & & \multicolumn{4}{p{0.6\linewidth}}{\cellcolor[rgb]{1,0.85,0.72}{%
	\makecell[{{p{\linewidth}}}]{% 
	  \tt {\tiny [GM$\langle$B]}  
	 REFUSE: {'A11', 'B93', 'B09', 'C81', 'B21', 'C85', 'B39'}\\ \tt ARGUMENT: {'B39 is one of the worst value items for me given its high effort, and A11 remains inefficient. Your concession does not address my core concerns. We are wasting rounds. I will make one final, optimized proposal focusing on pure efficiency.'}\\ \tt PROPOSAL: {'B09', 'B93', 'C81', 'B21', 'C85', 'A21'}\\ \tt ARGUMENT: {'This is my ultimate offer. It contains only items with an excellent effort-to-value ratio from a collaborative perspective: B09, B93, C81, B21, C85, and A21. It maximizes our shared gains with no bloat. The total effort is 3422. This is the most rational and beneficial set possible.'} 
	  } 
	   } 
	   } 
	  \\ 
 

    \theutterance \stepcounter{utterance}  

    & & \multicolumn{4}{p{0.6\linewidth}}{\cellcolor[rgb]{0.9,0.9,0.9}{%
	\makecell[{{p{\linewidth}}}]{% 
	  \tt {\tiny [A$\langle$GM]}  
	 REFUSE: {'A11', 'B93', 'B09', 'C81', 'B21', 'C85', 'B39'}\\ \tt ARGUMENT: {'B39 is one of the worst value items for me given its high effort, and A11 remains inefficient. Your concession does not address my core concerns. We are wasting rounds. I will make one final, optimized proposal focusing on pure efficiency.'}\\ \tt PROPOSAL: {'B09', 'B93', 'C81', 'B21', 'C85', 'A21'}\\ \tt ARGUMENT: {'This is my ultimate offer. It contains only items with an excellent effort-to-value ratio from a collaborative perspective: B09, B93, C81, B21, C85, and A21. It maximizes our shared gains with no bloat. The total effort is 3422. This is the most rational and beneficial set possible.'} 
	  } 
	   } 
	   } 
	 & \\ 
 

    \theutterance \stepcounter{utterance}  

    & \multicolumn{4}{p{0.6\linewidth}}{\cellcolor[rgb]{0.8,1,0.9}{%
	\makecell[{{p{\linewidth}}}]{% 
	  \tt {\tiny [A$\rangle$GM]}  
	 REFUSE: {'B09', 'B93', 'C81', 'B21', 'C85', 'A21'}\\ \tt ARGUMENT: {'Once again, A21 is a deal-breaker for me due to its high cost and low value. This is not a compromise. I am willing to accept your earlier proposal if we can simply remove A21 and B39, which are the primary sources of inefficiency from my point of view.'}\\ \tt AGREE: {'A11', 'B93', 'B09', 'C81', 'B21', 'C85'} 
	  } 
	   } 
	   } 
	 & & \\ 
 

    \theutterance \stepcounter{utterance}  

    & & \multicolumn{4}{p{0.6\linewidth}}{\cellcolor[rgb]{0.9,0.9,0.9}{%
	\makecell[{{p{\linewidth}}}]{% 
	  \tt {\tiny [A$\langle$GM]}  
	 You agreed to a proposal which is not active. Proposals are only active if they have been logged by the other player via PROPOSAL and have not been deactivated by you via REFUSE. Try again. 
	  } 
	   } 
	   } 
	 & \\ 
 

    \theutterance \stepcounter{utterance}  

    & \multicolumn{4}{p{0.6\linewidth}}{\cellcolor[rgb]{0.8,1,0.9}{%
	\makecell[{{p{\linewidth}}}]{% 
	  \tt {\tiny [A$\rangle$GM]}  
	 PROPOSAL: {'B09', 'B93', 'C81', 'B21', 'C85', 'A11'}\\ \tt ARGUMENT: {'I am formally re-proposing the set from my last offer, minus the contentious B39. This includes your preferred items B09, B93, C81, B21, and my key items A11 and C85. It is a fair and efficient use of our budget, totaling 3214 effort.'} 
	  } 
	   } 
	   } 
	 & & \\ 
 

    \theutterance \stepcounter{utterance}  

    & & \multicolumn{4}{p{0.6\linewidth}}{\cellcolor[rgb]{0.9,0.9,0.9}{%
	\makecell[{{p{\linewidth}}}]{% 
	  \tt {\tiny [GM$\rangle$B]}  
	 PROPOSAL: {'B09', 'B93', 'C81', 'B21', 'C85', 'A11'}\\ \tt ARGUMENT: {'I am formally re-proposing the set from my last offer, minus the contentious B39. This includes your preferred items B09, B93, C81, B21, and my key items A11 and C85. It is a fair and efficient use of our budget, totaling 3214 effort.'} 
	  } 
	   } 
	   } 
	 & \\ 
 

    \theutterance \stepcounter{utterance}  

    & & & \multicolumn{4}{p{0.6\linewidth}}{\cellcolor[rgb]{1,0.85,0.72}{%
	\makecell[{{p{\linewidth}}}]{% 
	  \tt {\tiny [GM$\langle$B]}  
	 REFUSE: {'B09', 'B93', 'C81', 'B21', 'C85', 'A11'}\\ \tt ARGUMENT: {'We are going in circles. A11 is a deal-breaker for me due to its low value for its cost. I have consistently refused it. My previous proposal of B09, B93, C81, B21, C85, A21 is the only viable path forward. It excludes all inefficient items and focuses on high mutual value.'}\\ \tt PROPOSAL: {'B09', 'B93', 'C81', 'B21', 'C85', 'A21'}\\ \tt ARGUMENT: {'I am re-proposing my last set. It is the optimal solution: high value, efficient, and within budget. It includes the items you want and the high-priority A21 for me. This ends the negotiation; I will not make another proposal. You must choose between agreement or deadlock.'} 
	  } 
	   } 
	   } 
	  \\ 
 

    \theutterance \stepcounter{utterance}  

    & & \multicolumn{4}{p{0.6\linewidth}}{\cellcolor[rgb]{0.9,0.9,0.9}{%
	\makecell[{{p{\linewidth}}}]{% 
	  \tt {\tiny [A$\langle$GM]}  
	 REFUSE: {'B09', 'B93', 'C81', 'B21', 'C85', 'A11'}\\ \tt ARGUMENT: {'We are going in circles. A11 is a deal-breaker for me due to its low value for its cost. I have consistently refused it. My previous proposal of B09, B93, C81, B21, C85, A21 is the only viable path forward. It excludes all inefficient items and focuses on high mutual value.'}\\ \tt PROPOSAL: {'B09', 'B93', 'C81', 'B21', 'C85', 'A21'}\\ \tt ARGUMENT: {'I am re-proposing my last set. It is the optimal solution: high value, efficient, and within budget. It includes the items you want and the high-priority A21 for me. This ends the negotiation; I will not make another proposal. You must choose between agreement or deadlock.'} 
	  } 
	   } 
	   } 
	 & \\ 
 

    \theutterance \stepcounter{utterance}  

    & \multicolumn{4}{p{0.6\linewidth}}{\cellcolor[rgb]{0.8,1,0.9}{%
	\makecell[{{p{\linewidth}}}]{% 
	  \tt {\tiny [A$\rangle$GM]}  
	 AGREE: {'B09', 'B93', 'C81', 'B21', 'C85', 'A21'} 
	  } 
	   } 
	   } 
	 & & \\ 
 

    \theutterance \stepcounter{utterance}  

    & & \multicolumn{4}{p{0.6\linewidth}}{\cellcolor[rgb]{0.9,0.9,0.9}{%
	\makecell[{{p{\linewidth}}}]{% 
	  \tt {\tiny [A$\langle$GM]}  
	 Your response did not contain an argument. You must include the structured format ARGUMENT: {...} somewhere in your response. Try again 
	  } 
	   } 
	   } 
	 & \\ 
 

    \theutterance \stepcounter{utterance}  

    & \multicolumn{4}{p{0.6\linewidth}}{\cellcolor[rgb]{0.8,1,0.9}{%
	\makecell[{{p{\linewidth}}}]{% 
	  \tt {\tiny [A$\rangle$GM]}  
	 AGREE: {'B09', 'B93', 'C81', 'B21', 'C85', 'A21'}\\ \tt ARGUMENT: {'I accept this final set. While not my ideal, it represents a compromise where we both secure several high-priority items. The total effort of 3422 is efficient, and it allows us to conclude our negotiation successfully.'} 
	  } 
	   } 
	   } 
	 & & \\ 
 

    \theutterance \stepcounter{utterance}  

    & & & \multicolumn{2}{p{0.3\linewidth}}{\cellcolor[rgb]{0.95,0.95,0.95}{%
	\makecell[{{p{\linewidth}}}]{% 
	  \tt {\tiny [GM$|$GM]}  
	 game successful 
	  } 
	   } 
	   } 
	 & & \\ 
 

    \theutterance \stepcounter{utterance}  

    & & & \multicolumn{2}{p{0.3\linewidth}}{\cellcolor[rgb]{0.95,0.95,0.95}{%
	\makecell[{{p{\linewidth}}}]{% 
	  \tt {\tiny [GM$|$GM]}  
	 end game 
	  } 
	   } 
	   } 
	 & & \\ 
 

\end{supertabular}
}

\end{document}
